% ***************************************************
% Executive Summary / Abstract
% ***************************************************

Marine ecosystems, particularly coral reefs, are undergoing unprecedented rates of degradation due to climate-driven thermal stress, bleaching events, and anthropogenic disturbances. Effective ecological conservation and reef restoration necessitate high-frequency, millimetre-accurate three-dimensional structural monitoring. However, traditional marine surveys depend heavily on diver-operated, two-dimensional planar transects and localized quadrats that fail to capture complex volumetric changes, structural breakage, and coral colony growth over time. 

This research project investigates the feasibility and efficacy of primitive-space Scene Change Detection (\gls{scd}) integrated with underwater-optimised 3D Gaussian Splatting (\gls{3dgs}) architectures. Unlike conventional render-then-compare paradigms that suffer from viewpoint dependencies and require pixel-space differencing, direct primitive-space change detection (\gls{oscd}) analyses structural transformations directly within native Gaussian attributes (position, anisotropic covariance, and colour harmonics).

To date, the project has achieved a major preliminary computational milestone: \textbf{15 out of 15 (100\%) full-lifecycle benchmark executions} have succeeded across an automated 6-stage pipeline on the QUT Aqua High-Performance Computing (\gls{hpc}) cluster. Across eight distinct state-of-the-art models evaluated on two real-world underwater shipwreck benchmarks (\textit{Submerged3D} Kwajalein and Tokai Maru) and a dual-state change detection dataset, the pipeline has extracted photorealistic 3D Gaussian representations, watertight surface meshes via Surface-Aligned Gaussian Splatting (\gls{sugar}), and high-fidelity novel view syntheses. Quantitative benchmarking confirms high reconstruction accuracy, with the primitive-space \gls{oscd} engine achieving a Peak Signal-to-Noise Ratio (\gls{psnr}) of $27.08\text{ dB}$, a Structural Similarity Index (\gls{ssim}) of $0.8315$, and a perceptual loss (\gls{lpips}) of $0.3440$.

Importantly, these preliminary executions serve as an essential proof-of-concept benchmark to rigorously evaluate candidate models. The ultimate objective of this thesis requires sensitive, high-density temporal photogrammetry from the \textbf{Australian Institute of Marine Science (\gls{aims}) EcoRRAP project} (a massive 1.35~TB dataset). This dataset is currently undergoing secure inter-institutional data transfer, subject to legal and data-sharing agreements. Building upon the empirical benchmarks established herein, Phase~2 and Phase~3 of the project will transplant the top-performing underwater 3DGS architectures (specifically \textit{3D-UIR} and \textit{SeaSplat}) into the \gls{oscd} pipeline, replacing naive terrestrial splatting backbones. These integrated models will undergo extended cluster optimization on Aqua to eliminate false-positive water column floaters and execute millimetre-precision primitive-space volumetric change tracking for marine conservation.
\vspace{1.5em}
\noindent\textbf{Keywords:} 3D Gaussian Splatting (3DGS), Underwater 3D Reconstruction, Primitive-Space Scene Change Detection, Surface-Aligned Regularization (SuGaR), Optical Attenuation, Marine Photogrammetry, Coral Reef Monitoring, EcoRRAP, High-Performance Computing (HPC).